\chapter{Introduction}
\label{chap:intro}

\section{Motivation}
\label{sec:motivation}
Retrieving items from a shopping list is a deceptively difficult robotics task.
It requires three capabilities at once: \emph{perception}, to see and identify
the correct object among clutter on a shelf; \emph{language understanding}, to
determine which item a natural-language request refers to; and
\emph{manipulation}, to physically grasp the object and place it accurately in a
receptacle. Traditionally, each of these is addressed by a separate, hand-engineered
subsystem, and integrating them is brittle and labour-intensive.

The task is also of practical interest. Service and retail robotics --- stocking
shelves, fulfilling orders, assisting shoppers --- is a rapidly growing application
area in which robots must operate in unstructured, human-centric environments and
respond to natural requests. A robot that can be told, in ordinary language, what to
fetch, and that can then perceive and manipulate everyday products, is a building
block for such applications. Studying this capability in simulation, under a tight
hardware budget, is a way to understand what is achievable with accessible resources
before committing to expensive hardware and data collection.

Vision--Language--Action (VLA) models offer an appealing alternative. By taking a
vision--language backbone that has already learned to perceive images and follow
language, and extending its output space from text to robot actions, a
\emph{single} network can map camera images and an instruction directly to motor
commands \parencite{zitkovich2023rt2, sapkota2026vla}. This promises generalist
robot control without a bespoke pipeline for every object and instruction.

The practical obstacle is scale. Capable VLAs typically contain billions of
parameters and are trained and deployed on large accelerators. For a student or a
low-cost robot, the relevant question is not ``what is the largest possible
model?'' but ``\emph{what can be fine-tuned and run on a single consumer GPU?}''
This thesis takes that constraint seriously --- a 12~GB GPU --- and asks whether a
compact VLA can nonetheless perform a useful, language-driven manipulation task
reliably.

The constraint is not merely a matter of convenience. It changes the design space:
it rules out the largest and most capable models, forces a choice about which parts
of a pretrained network to adapt, and makes data efficiency and evaluation
discipline decisive rather than optional. A recurring finding in this thesis, and in
the recent literature, is that under such constraints the quality of the data and of
the evaluation matters more than the raw capacity of the model. The work is
therefore as much about \emph{methodology} --- how demonstrations are collected, how
success is defined and measured, and how a deployment checkpoint is chosen --- as it
is about the model itself.

\section{Problem statement and scope}
\label{sec:problem}
The task studied in this thesis is a language-conditioned pick-and-place in a
simulated supermarket. A mobile manipulator --- a UR10e arm mounted on an Omron
mobile base, fitted with a Robotiq two-finger gripper --- is given a text
instruction naming a grocery item (for example, ``pick up the milk carton and
place it in the basket''). It must locate that item on a shelf, grasp it, and
place it in an onboard basket. The work is conducted entirely in the MuJoCo
physics simulator, which provides a safe, reproducible and low-cost setting for
benchmarking manipulation policies before they reach hardware
\parencite{todorov2012mujoco, nasiriany2024robocasa}. The base remains parked
during each pick; language-guided navigation of the store is identified as future
work (\refsec{sec:future}).

\section{Aim and objectives}
\label{sec:aims}
The aim of this thesis is to determine whether a compact, openly available VLA can
be fine-tuned and deployed, under a 12~GB GPU constraint, to perform reliable
language-conditioned grasping of everyday objects in simulation. The objectives
are:
\begin{enumerate}
    \item To build a reproducible simulated supermarket environment and a scripted
    expert capable of generating clean demonstrations of language-labelled
    pick-and-place (\refChap{chap:methodology}).
    \item To fine-tune the compact SmolVLA model on these demonstrations within the
    hardware budget, and to identify the training configuration that fits and
    generalises.
    \item To evaluate the resulting policy rigorously --- using closed-loop success
    rate, adequate sample sizes and confidence intervals --- and to select the
    deployed model by task success rather than training loss.
    \item To diagnose the dominant failure modes and to improve the policy through
    targeted, evidence-based changes.
    \item To assess whether the policy genuinely uses the language instruction, and
    to demonstrate an end-to-end system that collects a multi-item list.
\end{enumerate}

\section{Research questions}
\label{sec:rqs}
The objectives above are framed by four research questions that the thesis answers
empirically:
\begin{description}
    \item[RQ1 (Feasibility).] Can a compact, openly available VLA be fine-tuned and
    run on a single 12~GB GPU to perform language-conditioned pick-and-place at a
    useful success rate? Here ``useful'' is fixed in advance as a feasibility
    threshold --- the system completes the task more often than not, reliably enough
    to study --- and deliberately not as a deployment threshold, which for
    shelf-picking beside people would demand a rate far closer to unity than anything
    reported in this thesis. (\refsec{sec:results-headline},
    \refsec{sec:results-robustness})
    \item[RQ2 (Model vs.\ methodology).] At this data and hardware scale, do
    improvements come more from model capacity or from the quality of the data,
    metrics and evaluation? (\refsec{sec:results-journey}, \refsec{sec:results-diagnosis})
    \item[RQ3 (Grounding).] Does the fine-tuned policy genuinely condition on the
    language instruction, or does it exploit a fixed spatial correlate --- and if it
    does condition on language, is that conditioning implemented by localising the
    named object or by recalling where it usually is?
    (\refsec{sec:results-language}, \refsec{sec:results-lookup})
    \item[RQ4 (Failure structure).] When the policy fails, where in the
    perception--language--grasp--place chain does it fail, and can targeted,
    evidence-based changes correct it without increasing model size?
    (\refsec{sec:results-diagnosis})
\end{description}
These questions recur throughout, and are revisited together in the synthesis of
\refsec{sec:results-synthesis} and the conclusions of \refChap{chap:conclusions}.

\section{Contributions}
\label{sec:contributions}
The principal contributions of this thesis are:
\begin{itemize}
    \item A complete, reproducible pipeline --- scripted expert, demonstration
    collection, SmolVLA fine-tuning and closed-loop evaluation --- that operates
    within 12~GB of GPU memory by freezing the vision--language backbone and
    training only the action expert.
    \item A rigorous evaluation demonstrating an 80\% task-completion rate for the
    selected checkpoint (95\% Wilson confidence interval 70--87, $n=80$), drawn from
    320 closed-loop rollouts across four checkpoints, with the deployed checkpoint
    selected by task success rather than loss.
    \item A documented failure analysis that raised placement success from 0\%
    ($n=12$) to 80\% by diagnosing near-miss behaviour and addressing observability and
    episode-boundary issues, rather than by scaling the model or blindly
    retraining.
    \item Evidence that fixing the success metric and improving demonstration
    variety yielded larger gains than model choice at this data and hardware
    regime, echoing findings in the wider literature
    \parencite{black2025pi05, bharadhwaj2024roboagent}; supported by a same-task
    capacity control in which a 52-million-parameter policy with no language input
    matches the 450-million-parameter VLA's manipulation quality, so that the VLA's
    contribution is identified as item \emph{selection} rather than motor competence.
    \item A controlled characterisation of what language conditioning amounts to at
    this data scale: the policy is genuinely instruction-driven, but a displacement
    test over 216 trials shows it maps each item word to a memorised trajectory rather
    than localising the named object visually --- together with the consequence, that
    task completion falls from 80\% to 58.8\% when objects are displaced by
    \SI{2.5}{cm}. Alongside this, a light-weight interpretability view of the policy's
    predicted action plan.
    \item A functioning end-to-end orchestrator that collects a multi-item shopping
    list, quantified over 40 randomised lists: whole-list success is 45\% for two
    items and 30\% for three. The degradation is shown \emph{not} to be caused by the
    accumulating basket, the explanation previously assumed, since success does not
    fall as the basket fills; the mechanism is left open, with a candidate consistent
    with the displacement findings above.
\end{itemize}
Collectively, these contributions are less about a new model or algorithm than about
a disciplined account of what a compact VLA can and cannot do under a realistic
hardware constraint, and about the methodological choices --- in data, metrics and
evaluation --- that determine the difference between an apparent result and a real one.
The recurring finding, established across the pilot and the main study, is that at this
scale those choices dominate model capacity, which is a practically useful message for
anyone building robots without access to frontier compute.

\section{Thesis structure}
\label{sec:structure}
\refChap{chap:background} reviews VLA models, the specific architecture of
SmolVLA, the ACT baseline, and the simulation and benchmark context.
\refChap{chap:methodology} describes the simulation environment, the scripted
expert, demonstration collection, the fine-tuning procedure and the evaluation
methodology. \refChap{chap:results} presents the results, including the journey
from the RoboCasa benchmark to the deployed system, the headline and per-item
success rates, the failure analysis, the language-grounding test, checkpoint
selection and the end-to-end orchestrator. \refChap{chap:discussion} interprets
these results, states the limitations honestly and outlines the ideal
unconstrained system. \refChap{chap:conclusions} concludes and describes future
work.
